\begin{table}[!htbp]
\centering
\begin{threeparttable}
\caption{Headline conclusions under randomised parameters and departures from the assumptions.}
\label{tab:robust-design}
\small\setlength{\tabcolsep}{4pt}
\begin{tabular}{lcccc}
\toprule
 & Random & Heavy tails & Outliers & Tested seas. \\
\midrule
TimesFM-3 cell wins (of 36) & 14 & 15 & 16 & 16 \\
TimesFM-3 median ratio to cell-best & 1.029 & 1.038 & 1.004 & 1.017 \\
TimesFM-3 worst ratio to cell-best & 1.56 & 1.54 & 1.49 & 1.55 \\
Smallest worst ratio, classical & 2.12 & 2.32 & 2.09 & 1.58 \\
Significant comparisons, won--lost & 93--29 & 100--28 & 101--27 & 100--27 \\
D4, $n \geq 48$: AutoARIMA better (of 3) & 3 & 3 & 2 & 3 \\
D8: Croston better on RMSSE (of 4) & 4 & 4 & 4 & 4 \\
D5, $n = 24$: AutoETS / seasonal naive ($m = 12$) & 2.03 & 2.05 & 1.91 & 2.25 \\
\bottomrule
\end{tabular}
\begin{tablenotes}[flushleft]\footnotesize
\item \textit{Note:} 200 replications per (process, length) cell, 7\,200 series per column; every series draws its own parameters from the ranges in Table~S8 of the Supporting Information. Ratios use mean MASE over $h = 1, \dots, 12$. Significant comparisons are paired Wilcoxon tests of TimesFM-3 against each of the five general-purpose classical methods in each cell (180 per column), Benjamini--Hochberg at 0.05 within each (column, opponent) family of 36 cells (the main design's families also include the three horizon slices). The D8 row compares TimesFM-3 with the best of Croston-SBA, TSB and ADIDA; the last row divides by seasonal naive with the true period in every column. In the tested-seasonality column the classical methods use $m = 12$ only where the M4 seasonality test finds seasonality; TimesFM-3's forecasts are those of the first column.
\end{tablenotes}
\end{threeparttable}
\end{table}
